% Ad Familiares 13.15 --- headnote
% ad-familiares-13-15, end of May 45 BC, Rome

\textit{Cicero at Rome to Gaius Julius Caesar in Spain --- or, by
the time the letter overtook him, on the road back from the
Munda campaign --- written at the very end of May 45 BC (the
manuscript dateline: \textit{Scr. Romae ex. m. Mai. a. 709
(45)}). The salutation \textit{CICERO CAESARI IMP. S.} is the
formula Cicero uses to Caesar in his \textit{imperator}-capacity
after Pharsalus and during the dictatorship. This is one of the
small number of surviving private letters from Cicero to Caesar
himself, and the only one of \textit{Familiares} 13's
seventy-nine commendations addressed to him --- which is part of
what makes it so peculiar.}

\textit{The ostensible business is brief: a recommendation of
the younger Precilius, son of an old friend, on the standard
courtesy-and-clientship terms. Around that nucleus Cicero builds
an extraordinary frame of Homeric and Euripidean quotation,
seven Greek tags in three sections --- a density unmatched
anywhere else in the \textit{Familiares}. The line of thought
runs: the boy's father used to scold me for not joining your
side; I refused (Iliad 9.587, Phoenix on Meleager); the great
men kept urging me to glory (Odyssey 1.302, on being valorous
for the sake of posterity); a black cloud of grief covered me
(Iliad 17.591); but now I have moved on from Homer's grandiloquence
to Euripides' realism, ``I hate the wise man who is not wise on
his own account.'' The pose is self-mockery at the Pompeian
choice, and a flattery of Caesar's victory: the speaker has
learned, late, the lesson Caesar's friends had been trying to
teach him. Cicero closes by drawing the reader's attention to
the device: \textit{genere novo sum litterarum ad te usus, ut
intellegeres non vulgarem esse commendationem} --- ``I have
employed with you a novel kind of letter, that you might
understand this is no ordinary recommendation.''}

\textit{The piece is a delicate one. Cicero is a defeated
Pompeian writing to the victor of the civil war, asking a
personal favour, and folding into the request a half-apology
for his own past resistance to that victory --- carried, in
Greek, at one remove from his own voice. The cluster of
Homeric tags is unmistakably the literary virtuosity of an
ageing consular displaying that he can still play the game;
the Euripidean turn at the end carries an edge, since the
``wise man who is not wise on his own behalf'' may be Cicero
himself in his Pompeian years, but may also stand at a glancing
angle to the addressee. The metadata entry carries a year-
precision placeholder of \texttt{-0045-08-28}; the Perseus
dateline points to the end of May 45, which should be carried
into the meta entry at the next consolidation pass.}
